% 2011 年安徽卷文科第 10 题：函数 f(x)=a x^n (1-x)^2 在 [0,1] 上的图象，判断 n 的可能值。
% 来源 PDF：普通高考/2011/2011安徽文.pdf，SHA-256 af8382bc2cd47a1d5e08bb622275f77ad100aa4f844b5cfd2fd5cd217de58ffc。
% 原页：第 1 页（A3 三栏）第 2 栏中部，书页码 645。
% 重绘依据：照原页——只画第一象限的直角坐标系，x 轴标 0.5 与 1，y 轴标 0.5（与曲线最高点同高），
%   曲线自原点 O 上升、在 0.5 左侧取最高点、于 x=1 处与 x 轴相切后结束；
%   曲线按原页形状（峰在 0.5 左侧、峰高 0.5、右端水平）取 n=1、a=27/8 的函数绘制，只取形状，不标函数式。
\begin{tikzpicture}[scale=3.1, >=stealth, semithick]
  \draw[->] (-0.15,0) -- (1.26,0) node[right] {$x$};
  \draw[->] (0,-0.08) -- (0,0.66) node[above] {$y$};
  \draw[domain=0:1, samples=240, smooth] plot (\x, {3.375*\x*(1-\x)*(1-\x)});
  \draw (0.5,0) -- (0.5,0.013);
  \node[below=0.04cm] at (0.5,0) {$0.5$};
  \draw (1,0) -- (1,0.013);
  \node[below=0.04cm] at (1,0) {$1$};
  \draw (0,0.5) -- (-0.013,0.5);
  \node[left=0.06cm] at (-0.013,0.5) {$0.5$};
  \node[below left=0.01cm and 0.01cm] at (0,0) {$O$};
\end{tikzpicture}
